\section{Boundary and conclusion}

Deterministic regular receipt lifts of a finite directed cycle have a complete
and economical classification. Local edge voltages reduce under gauge to one
full-circuit holonomy conjugacy class. The order of that receipt fixes the
entire cycle profile:
\[
\text{component length}=n\,\ord(h),
\qquad
\text{component count}=\frac{|R|}{\ord(h)}.
\]

Connectedness is maximally restrictive. It requires the retained receipt to
generate the entire fiber group and therefore requires that group to be
cyclic. Nonabelian receipt groups still admit regular lifts, but no single
holonomy can connect all receipt states over one visible cycle.

The theorem does not derive a receipt group or select its holonomy from a
physical system. Receipt order is not identified with physical time, distance,
energy, winding, or action magnitude. Such a representation would have to
construct the finite group action and circuit receipt independently.

The earned statement is:
\begin{quote}
A receipt does not merely remember the circuit. Its order determines how many
circuits are required for complete return.
\end{quote}

